\chapter*{Preface}
A sphere, a doughnut and a coffee cup are shapes, and so is the set of all rotations of space, the set of all unitary matrices, or the
surface on which modular forms live. This book is about such shapes: how to do calculus on them, how to measure their curvature, how to
recognise that two of them are different, and how to integrate over them. Its hero is the idea, due to Riemann, that a space can be
curved without sitting inside anything else, and that everything one needs to know about it can be read from the inside.

The first four chapters are differential geometry. Smooth manifolds, built from the inverse function theorem; vector fields and the
flows they generate, with Lie's bracket and Frobenius' theorem; differential forms, which turn the theorems of Green, Gauss and Stokes
into a single line and prove Brouwer's fixed point theorem in every dimension; and curvature, from Gauss's Theorema Egregium to the
Gauss--Bonnet theorem, where a local quantity integrates to a global invariant.

The next two chapters are algebraic topology. The fundamental group and covering spaces, where the circle's winding number becomes a
group and the subgroups of free groups turn out to be free; then homology, which counts holes in every dimension and proves that
open sets of $\R^n$ cannot be open sets of $\R^m$ unless $n=m$.

The last five chapters bring geometry back to the subjects of this series. Matrix Lie groups, where the rotations of space are
quaternions in disguise and a full turn is not the same as no turn at all; Haar measure and Weyl's integration formula, which give
the law of the eigenvalues of a random unitary matrix, the bridge to random matrices and the zeros of the zeta function; the
hyperbolic plane, the modular group and the geodesic flow, in which continued fractions become the records of geodesics crossing a tiling; the Riemann mapping theorem, which sends every simply connected domain other than the plane onto the disc, after the half-plane has been written onto the disc by an explicit map; and integral
geometry, where Buffon's needle grows into Crofton's formula and the measure of all lines in the plane.

Eleven appendices, A to K, hold the elementary and pictorial geometry: transformations and Klein's programme, inversion, projective
geometry, conics and quadrics, plane curves and cubics, lattices and sphere packings, non-Euclidean geometry, convex bodies, and the
curvature of curves and polyhedra. Each opens by saying what it connects and what it assumes, and the chapters cite them where they need them.

The book assumes Volume~7 (topology, measure and integration), the complex analysis of Volume~5 and the homology form of Cauchy's theorem in
Volume~8, the Haar measure of Volume~8, the groups and quaternions of Volume~3, the fundamental domain of the modular group in Volume~6, and the
continued fractions of Volume~2. Each chapter ends with exercises, each with a hint and a complete solution.

The series asks whether pure order and pure chance can be told apart, and why the primes look random yet obey laws. In this volume two theorems carry
that question. Gauss--Bonnet shows that curvature, measured point by point from inside a surface, adds up to a number that counts its holes. Weyl's
formula shows that the eigenvalues of a random unitary matrix repel each other with an exact law, the law Montgomery found among the zeros of the zeta
function. Volume~10 takes up the structures common to all of these objects. The group law on a cubic curve, begun in Appendix~G, is completed in Volume~13.
